\chapter{Storage: the schema, batching, and the log counter}\label{ch:storage}

\texttt{src/ate/storage/} is the only place ATE writes SQL. Three files: \texttt{schema.py} declares
every table as a SQLAlchemy Core table, \texttt{store.py} implements \texttt{RunStore}, the only
interface the rest of the application uses, and \texttt{backup.py} holds the one routine that backs a
database up and then migrates it. Nothing outside these files contains a table name or a query.

\section{The record, and the version}

The Test Run record is ten tables. \texttt{tests/\allowbreak{}test\_\allowbreak{}manuals.py} reads them from
\texttt{schema.metadata.tables}, so this list is the code's list:

\begingroup\small
\begin{longtable}{p{3.2cm} p{8.4cm}}
\toprule
\textbf{Table} & \textbf{One row is} \\
\midrule
\endhead
\texttt{test\_run} & one Test Run: Operator, Project, UUT, schedule, Media Release, Stop on Fail, start
  and end times, UTC offset, Result, End Reason \\
\texttt{node} & one step or header in the expanded tree, with its number, enable state, Result and
  Attempt count \\
\texttt{attempt} & one execution of one step, with its times, Result, message, actual, expected, and the
  Stop-on-Fail choice made at it \\
\texttt{action} & one Pre-test or Post-test Action call, with phase, order, times and outcome
  (\texttt{ok}, \texttt{error}, \texttt{aborted}, \texttt{skipped}) \\
\texttt{log\_line} & one line of the Run log, owned by an Action or an Attempt, or run-level \\
\texttt{input\_value} & one Step Input answer \\
\texttt{equipment} & one \texttt{ATE.addEquipment} declaration, in declaration order \\
\texttt{media\_release} & one Media Release in the resolved chain; the release owning the run is seq 0 \\
\texttt{run\_file} & one file the run used, with expected and computed checksums and whether they differed \\
\texttt{schedule\_instance} & the top-level schedule and every Sub-schedule include, in tree order \\
\bottomrule
\end{longtable}
\endgroup

Types are deliberately portable: UUIDs are 36-character text, times are ISO-8601 UTC text with
milliseconds (\texttt{\_now\_iso}), booleans are 0/1 (\texttt{\_as\_int}), structured values are JSON
text. There is no SQLite-only column type anywhere; a Postgres backend is meant to be a connection
change. The version lives in SQLite's \texttt{PRAGMA user\_version}, and \texttt{SCHEMA\_VERSION} is
currently 5.

At launch, \texttt{check\_\allowbreak{}database\_\allowbreak{}at\_\allowbreak{}launch} inspects the file read-only --- it never creates it ---
and compares \texttt{user\_version} with \texttt{SCHEMA\_VERSION}. A missing, empty or older database
refuses with \texttt{run `ate db migrate` to create or upgrade the database at \dots}; a newer one refuses
with \texttt{this ATE is too old for the database at \dots; update ATE}. Only an exact match allows
launch. The application never migrates; \texttt{ate db migrate} does, and the hint carries
\texttt{--config <path>} when the Station is not reading the platform default (packet 30c).

\texttt{RunStore.\_\allowbreak{}ensure\_\allowbreak{}schema} is idempotent: \texttt{metadata.create\_all} reflects each table through
the backend's own catalog, so reopening a populated database never fails. The version stamp is what
packet 35 changed: \texttt{PRAGMA user\_version} is raised only when the build's
\texttt{missing\_tables(path)} was non-empty --- when a table was really created. A database that stores
an older version but already holds every table this ATE declares has had no migration to earn the stamp,
so opening it raises \texttt{SchemaVersionError} naming that version instead of quietly declaring it
current. WAL and the version are the two pragmas set on the store's own stdlib connection: the only raw
SQL in ATE, and it is confined to pragmas.

\texttt{src/ate/storage/backup.py} is the one routine \texttt{ate db migrate} and \texttt{ate setup}
both call, so an upgrade through an installer and an upgrade by hand take the same copy.
\texttt{backup\_for\_migration} decides honestly whether a copy is needed: none for a database that does
not exist, none for one with no tables, none for one that needs no migration; otherwise SQLite's
\texttt{Connection.backup} writes \texttt{<db>.backup-<YYYYMMDD-HHMMSS>} beside the file, never over an
existing name. The backup API is used rather than \texttt{cp} because it folds \texttt{-wal} into the
copy: a database with a Test Run in flight is copied whole. \texttt{migrate\_database} then refuses if
the copy cannot be written, and reports the copy's path when a migration fails after it.

\section{One commit per row, and why that is a design decision}

The Core Engine is autocommit: each write commits immediately. That is what makes the record
append-only and crash-safe --- an open run simply has no \texttt{ended\_at} yet --- and it is also the
reason an abort over 400 unrun steps used to make 404 commits and take longer than spec story 83's 10 s with
no stall at all (packet 32b). The answer was batching, not a different transaction model.

SQLite limits the number of variables one statement may carry; the lowest default is 999
(\texttt{\_\allowbreak{}SQL\_\allowbreak{}MAX\_\allowbreak{}VARIABLES}, the pre-3.32 default, so the code stays safe on old builds too).
\texttt{\_\allowbreak{}rows\_\allowbreak{}per\_\allowbreak{}statement(table)} computes the batch size from the table's column count ---
\texttt{999 /\allowbreak/\allowbreak len(table.columns)} --- rather than a hard-coded 80 chosen for \texttt{node}'s twelve
columns. A batch sized for one table breaks the day another table gains a column, which
\texttt{test\_\allowbreak{}a\_\allowbreak{}row\_\allowbreak{}batch\_\allowbreak{}is\_\allowbreak{}sized\_\allowbreak{}to\_\allowbreak{}the\_\allowbreak{}table\_\allowbreak{}it\_\allowbreak{}writes} is there to prevent.
\texttt{\_row\_chunks} then yields the rows in pieces that fit.

There is a trap here that a new writer will hit. Each chunk is its own commit, so the run's counters must
advance \emph{inside} the loop, only for chunks that landed. \texttt{insert\_log\_lines} does exactly
that: before packet 32b a half-written batch left \texttt{\_log\_seq} behind the rows already in the
table, and the next write failed with \texttt{UNIQUE constraint failed: log\_line.run\_uuid,
log\_line.seq}. The test is
\texttt{test\_\allowbreak{}a\_\allowbreak{}log\_\allowbreak{}batch\_\allowbreak{}failing\_\allowbreak{}part\_\allowbreak{}way\_\allowbreak{}through\_\allowbreak{}keeps\_\allowbreak{}the\_\allowbreak{}log\_\allowbreak{}counter\_\allowbreak{}true}.

\section{\texttt{\_log\_seq}: one counter per run, four tables}

\texttt{log\_line}, \texttt{action}, \texttt{input\_value} and \texttt{equipment} all have
\texttt{(run\_uuid, seq)} as their primary key, and they share one per-run counter held in
\texttt{RunStore.\_log\_seq}. That is what makes a log line, an Action and an Equipment declaration
uniquely keyed without a global sequence, and it is why the ordering of the drain calls matters: the
engine diffs the host's collected lines around each step, and a batch can hold several lines from one
Attempt.

Rules for a new row type that shares this key:

\begin{itemize}
  \item take \texttt{next\_\allowbreak{}seq = self.\_\allowbreak{}log\_\allowbreak{}seq.get(run\_\allowbreak{}uuid, 0)} and advance it by the number of rows
    that actually landed;
  \item never assume a batch is one commit;
  \item stamp ownership (\texttt{action\_seq}, or \texttt{node\_seq} plus \texttt{attempt\_no}) the way the
    log lines do, so the report can group a row under the step or Action that produced it.
\end{itemize}

\section{Reading, and the \texttt{\_run} subtlety}

Every read goes through \texttt{RunStore.\_run}, which executes a Core statement and
\texttt{fetchall()}s the rows \emph{inside} the connection context. Packet 25a review found why:
returning the open \texttt{Result} left a read statement alive after the context closed; the GUI thread
fetched it later, the connection pool handed the same sqlite3 connection --- cursor still open --- to the
run worker, and a worker INSERT colliding with the unfinalised read aborted a long run with ``database is
locked''. Do not return a result object from \texttt{\_run}. Materialise it.

\texttt{query\_runs} is the History screen's query: filters on the eight \texttt{HISTORY\_COLUMNS},
sorting on exactly those keys, and \texttt{run\_uuid} ordered last so paging never repeats or skips a row
when several runs share a sort value. Dates are normalised to ISO-8601 UTC by \texttt{\_norm\_date} so
lexical comparison is correct.

\cando{add a table to the record the way the code does it, and explain why a batch is not one commit.}
